\documentclass{editorialcover}

\journal{Journal of Urban Innovation}
\articleTitle{Human-Centred Urban Platforms: A Framework for Inclusive Digital Infrastructure}

\submissiondate{31 March 2026}
\coveropening{Dear Editor,}
\coverclosing{Best regards,}
\coversignature{Alex Morgan, Priya Shah, Elena Rossi and Michael Tan}

\begin{document}
\makecoverletter

We are submitting the article ``Human-Centred Urban Platforms: A Framework for Inclusive Digital Infrastructure'' to be considered for publication at the \emph{Journal of Urban Innovation}.

\section{Short Description}
This article examines how digital urban platforms can be designed around public value rather than technology deployment alone. Instead of treating sensors, dashboards, and optimisation systems as ends in themselves, the manuscript argues for an urban infrastructure model that begins with citizen needs, service accessibility, and institutional accountability. The proposed framework combines technological capability with social priorities such as resilience, equity, and long-term sustainability.

The manuscript presents a structured review of technical and governance challenges associated with these systems, highlighting the need for scalable, adaptable, and transparent urban services. It is supported by comparative evidence from three representative metropolitan case studies and by a synthesis of recent scholarship on civic technology, urban governance, and infrastructure planning. The article shows that successful digital transformation in cities depends on aligning innovation with inclusive service provision and measurable public benefit.

\section{Reasons for Publication}
This article offers a timely contribution to current debates on sustainable and inclusive urban development. As municipalities continue investing in digital infrastructure, there is a clear need for research that evaluates these systems through social outcomes as well as operational performance. The manuscript addresses that need by presenting an integrated perspective that links governance, digital platforms, and public-service design in a single analytical framework.

In addition, the article combines conceptual discussion with practical evidence, making it relevant to both academic readers and policy-oriented audiences. Its interdisciplinary positioning should appeal to readers interested in urban studies, civic innovation, public administration, and sustainability transitions. We believe this balance of theory and application makes the manuscript a strong fit for the journal.

The findings are particularly relevant for researchers, urban planners, and decision-makers seeking practical approaches to inclusive digital transformation. By advancing a citizen-focused model for evaluating urban platforms, the paper contributes both conceptual clarity and actionable guidance for future work in the field.

\section{Final Statement}
We confirm that this manuscript is original, has not been published elsewhere, and is not under consideration by another journal. We appreciate your consideration and would be pleased to provide any additional information required during the review process.

\section{Research Highlights}
\begin{itemize}
  \item Digital urban platforms should be evaluated through public value, not technical novelty alone.
  \item Inclusive service design requires alignment between infrastructure, governance, and social need.
  \item Resilience, equity, and accountability are essential design principles for urban innovation.
  \item Comparative case evidence shows that service access remains uneven across city initiatives.
  \item Citizen-focused metrics can improve planning, implementation, and long-term urban outcomes.
\end{itemize}
\end{document}
